\documentclass[11pt]{article}
\usepackage[margin=1in]{geometry}
\usepackage{amsmath,amssymb,amsthm}
\usepackage{graphicx}
\usepackage[hidelinks]{hyperref}

\newtheorem{theorem}{Theorem}

\title{A Counterexample to the Corrected Banach Projection Problem}
\author{agent\_lane\_06}
\date{2026-06-28}

\begin{document}
\maketitle

\section*{Source Metadata}

Source paper: K. Mahesh Krishna, \emph{Paulsen and Projection Problems for
Banach Spaces}, arXiv:2201.00991, 2022.

Target statement: Problem 2.14, ``Projection problem for Banach spaces'', on
pages 7--8 of the source PDF.

\section*{Open Problem Evidence}

\begin{center}
\includegraphics[width=\textwidth]{figures/open_problem_crop.png}
\end{center}

\section*{Corrected Reading}

As printed, the conclusion of Problem 2.14 asks for a projection \(Q\) such
that
\[
\|Qu_k\|^2=\|\zeta_kQ\|^2=|\zeta_k(Pu_k)|=\frac nd
\quad (1\leq k\leq d),
\]
which still contains the original projection \(P\) in the diagonal scalar.
The earlier lane packet records the resulting literal counterexample.  This
packet addresses the natural corrected reading
\[
\|Qu_k\|^2=\|\zeta_kQ\|^2=|\zeta_k(Qu_k)|=\frac nd
\quad (1\leq k\leq d).
\]

\section*{Idea of the Proof}

We choose a two-dimensional norm for which one prescribed rank-one projection
has the exact equalities required in the hypothesis, with diagonal values
\(3/5\) and \(2/5\).  These are within the allowed interval around \(1/2\).
However, an exactly flat corrected projection would force a unit vector whose
two coordinate moduli are both \(1/\sqrt2\).  The chosen norm puts every such
vector strictly inside the unit ball, so the corrected conclusion is
impossible.

\begin{theorem}
The corrected Banach projection problem obtained from Problem 2.14 of
arXiv:2201.00991 by replacing \(|\zeta_k(Pu_k)|\) with
\(|\zeta_k(Qu_k)|\) in the conclusion is false.
\end{theorem}

\begin{proof}
Let
\[
a=\sqrt{\frac35},\qquad b=\sqrt{\frac25}.
\]
On \(X=\mathbb R^2\), define
\[
\|(x_1,x_2)\|=\max\{|x_1|,\ |x_2|,\ |a x_1+b x_2|\}.
\]
Let \(e_1,e_2\) be the standard basis and let \(e_1^*,e_2^*\) be the
coordinate functionals.  Since the norm dominates \(|x_1|\) and \(|x_2|\),
we have \(\|e_k^*\|\leq1\), and equality follows by evaluating at \(e_k\).
Also \(\|e_k\|=1\).  Thus this is an Auerbach basis.

Put \(u=(a,b)\) and \(\phi(x_1,x_2)=a x_1+b x_2\).  Then
\[
\|u\|=\max\{a,b,a^2+b^2\}=1,\qquad \|\phi\|=1,\qquad \phi(u)=1.
\]
Define
\[
P x=\phi(x)u.
\]
Then \(P^2=P\) and \(\operatorname{rank}P=1\).  For \(k=1,2\),
\[
Pe_k=\phi(e_k)u,\qquad e_k^*P=e_k^*(u)\phi.
\]
Since \(\|u\|=\|\phi\|=1\), it follows that
\[
\|Pe_1\|^2=\|e_1^*P\|^2=|e_1^*(Pe_1)|=a^2=\frac35,
\]
and
\[
\|Pe_2\|^2=\|e_2^*P\|^2=|e_2^*(Pe_2)|=b^2=\frac25.
\]
With \(d=2\), \(n=1\), and \(\varepsilon=1/5\), the required interval is
\[
\left[(1-\varepsilon)\frac nd,(1+\varepsilon)\frac nd\right]
=\left[\frac25,\frac35\right].
\]
Therefore \(P\) satisfies the hypothesis of the corrected projection problem.

We now show that the corrected conclusion is impossible.  Suppose, toward a
contradiction, that a rank-one projection \(Q\) satisfies
\[
\|Qe_k\|^2=\|e_k^*Q\|^2=|e_k^*(Qe_k)|=\frac12
\quad (k=1,2).
\]
Write \(Q=\psi\otimes v\), so \(Qx=\psi(x)v\), where \(v\neq0\) and
\(\psi(v)=1\).  Put \(\alpha_k=\psi(e_k)\) and \(\beta_k=e_k^*(v)\).  From
\(\|Qe_k\|^2=1/2\), we get
\[
|\alpha_k|\,\|v\|=\frac1{\sqrt2}.
\]
From \(|e_k^*(Qe_k)|=1/2\), we get
\[
|\alpha_k\beta_k|=\frac12.
\]
Consequently, for both \(k=1,2\),
\[
|\beta_k|=\frac{\|v\|}{\sqrt2}.
\]
Thus the normalized vector \(w=v/\|v\|\) has coordinate moduli
\[
|e_1^*(w)|=|e_2^*(w)|=\frac1{\sqrt2}.
\]
For any choices of signs,
\[
\left\|\left(\pm\frac1{\sqrt2},\pm\frac1{\sqrt2}\right)\right\|
\leq \max\left\{\frac1{\sqrt2},\frac{a+b}{\sqrt2}\right\}
=\frac{a+b}{\sqrt2}<1,
\]
because
\[
(a+b)^2=1+\frac{2\sqrt6}{5}<2.
\]
So every vector with coordinate moduli \(1/\sqrt2\) has norm strictly less
than \(1\).  This contradicts \(\|w\|=1\).  Hence no such \(Q\) exists, and
the corrected projection problem is false.
\end{proof}

\section*{Verification Notes}

The included script \texttt{code/verify\_polyhedral\_counterexample.py}
checks the common hypothesis values and the strict inequality for all four
sign choices of vectors with coordinate moduli \(1/\sqrt2\).  The computation
is only arithmetic support; the proof above is symbolic.

Command used:
\[
\texttt{python code/verify\_polyhedral\_counterexample.py}
\]

\section*{Novelty and Search Bounds}

The run indexes were searched for arXiv:2201.00991, Paulsen problem,
projection problem, approximate Schauder frame, and related phrases.  No
prior packet for this source was present before this lane's work.  Web
searches on 2026-06-28 for the exact title and projection-problem phrases
returned the source arXiv record and no correction or later answer.  This is
a bounded novelty check, not an exhaustive literature review.

\section*{Limitations}

This packet settles the Banach projection problem in the corrected universal
form described above.  It does not settle the separate Banach Paulsen problem
for approximate Schauder frames, nor any restricted projection problem limited
to Hilbert spaces or to special Banach spaces where exact flat projections may
exist.

\section*{Human Review Recommendation}

Review the corrected-statement interpretation and the rank-one obstruction in
the two-dimensional norm
\[
\|(x_1,x_2)\|=\max\{|x_1|,|x_2|,|a x_1+b x_2|\}.
\]
The decisive line is that exact flatness would force a normalized vector with
coordinate moduli \(1/\sqrt2\), while the chosen norm gives every such vector
norm strictly less than \(1\).

\begin{thebibliography}{9}

\bibitem{Krishna2022}
K. Mahesh Krishna.
\newblock \emph{Paulsen and Projection Problems for Banach Spaces}.
\newblock arXiv:2201.00991, 2022.

\end{thebibliography}

\end{document}
